only finitely many singular points, a reduction justified by the projection formula X 4.4.4 and by 4.5 (ii). There therefore exists a closed integral curve @@M0@@ on @@M1@@ such that @@M2@@, and @@M3@@ is locally principal outside a finite set @@M4@@. Let @@M5@@ be the blowup of the ideal @@M6@@ of @@M7@@. The ideal @@M8@@ is invertible, and defines a subscheme of @@M9@@ which is the union of two subschemes @@M10@@, @@M11@@ such that @@M12@@ and @@M13@@. By the projection formula X 4.4.4, we have that @@M14@@ equals @@M15@@, and we may replace @@M16@@ by @@M17@@; we are therefore reduced to the preceding case, and the proof of 7.1 is complete.

\sgaStableTarget{sga6.e13.rem.7-2}\textbf{Remark 7.2.} It is unknown whether an irreducible proper surface (over an algebraically closed field), which possesses an invertible sheaf @@M18@@ such that @@M19@@, is necessarily projective. It is also unknown whether, for every irreducible proper surface, the intersection form on @@M20@@ is nondegenerate; it follows from 7.1 that, when it is degenerate, it is necessarily negative (nondefinite).

\sgaStableTarget{sga6.e13.cor.7-3}\textbf{Corollary 7.3.} \emph{Under the conditions of 7.1, let @@M21@@ be an invertible sheaf on @@M22@@ which is not numerically proportional to @@M23@@, and let @@M24@@ and @@M25@@. Then, @@M26@@.}

Indeed, put @@M27@@. Since @@M28@@, by 7.1, @@M29@@; whence the assertion.

